\makeatletter
\def\input@path{{./}{../}{../../}{../../../}{../../../../}{preamble/}{../preamble/}{../../preamble/}{../../../preamble/}}
\makeatother

\input{preamble_exams.tex}
\usepackage{needspace}
\usepackage{titling}
\setlength{\droptitle}{-2em}
\pretitle{\par\begingroup\centering\LARGE}
\posttitle{\par\endgroup}
\preauthor{}\postauthor{}
\predate{}\postdate{}
\maketitlehookd{\vspace{-1em}}
\linespread{1.12}
\setlength{\parskip}{5pt}
\newgeometry{margin=0.65in,top=0.55in,bottom=0.65in}

\title{Grade 11 Term 2 Home Practice -- C7\\\large Confidence and precision\\[0.25em]
	\normalsize Global Economics -- 2026--2027}
\author{}
\date{}

\begin{document}
	\maketitle

	\setcounter{tocdepth}{1}
	\setcounter{secnumdepth}{2}
	\phantomsection
	\label{toc}
	\tableofcontents

	\section*{Evaluated criteria}
	\begin{enumerate}[label=\textbf{C\arabic*},start=7]
		\item Compares confidence intervals constructed from samples of different sizes and confidence levels, analysing the trade-off between confidence and precision.
	\end{enumerate}

	\textbf{How to practise.} Try each task before opening its blue solution link. The arrows return to the contents, problem list or current problem.

	\section{Worked Solution Examples}
	\label{sec:worked-solution-examples}

	\SectionProblemsTableOfContents{
		\SectionProblemLink{prob:c7-1}{Mobile ticketing and sample size (C7)} \\
		\SectionProblemLink{prob:c7-2}{Mobile ticketing and the confidence target (C7)}
	}

	% ----------------------------------------------------------------
	\subsection{Mobile ticketing and sample size (C7)}
	\label{prob:c7-1}

	A transit authority compares simple random samples of 100 and 400 riders from the same registry of 20,000 eligible riders. Each sample is under 10\% of the registry. In each, 40\% report buying a mobile ticket last week. The counts for both outcomes are at least 10 in each sample. Use $z^*=1.645$ for 90\% confidence.

	\QuestionsTasks[sec:worked-solution-examples]{
		{7}/{sol:c7-1-a}/
		{Use $\widehat{SE}=\sqrt{\hat p(1-\hat p)/n}$ to construct and compare both 90\% intervals. Give each margin of error in percentage points.},
		{7}/{sol:c7-1-b}/
		{Use the interval widths to explain what the larger sample buys the transit authority at the same confidence level.}
	}

	\clearpage

	\Needspace{5\baselineskip}
	\phantomsection
	\label{sol:c7-1-a}
	\subsubsection*{C7 -- Construct the intervals at the same confidence}

	\noindent
	\begin{minipage}[t]{0.45\textwidth}
	\vspace{0pt}
	\centering
	\begin{tikzpicture}
	\begin{axis}[width=0.97\linewidth,height=4.3cm,axis lines=left,
	xmin=29,xmax=51,ymin=0.5,ymax=2.5,ytick={1,2},
	yticklabels={$n=100$,$n=400$},xtick={30,35,40,45,50},
	xlabel={Mobile-ticket use (\%)}]
	\addplot[black,dashed] coordinates {(40,0.5) (40,2.5)};
	\addplot[blue!65!black,thick,mark=|] coordinates {(31.9412,1) (48.0588,1)};
	\addplot[blue!65!black,thick,mark=|] coordinates {(35.9706,2) (44.0294,2)};
	\end{axis}
	\end{tikzpicture}
	\end{minipage}
	\hfill
	\begin{minipage}[t]{0.51\textwidth}
	\vspace{0pt}
	\centering
	\begin{tabular}{@{}rccc@{}}
	\toprule
	$n$ & $\widehat{SE}$ & $E$ (pp) & 90\% interval \\
	\midrule
	100 & 0.04899 & 8.06 & $[31.94\%,48.06\%]$ \\
	400 & 0.02449 & 4.03 & $[35.97\%,44.03\%]$ \\
	\bottomrule
	\end{tabular}\par\smallskip

	\raggedright
	Both centres are $40\%$. For 100 riders, $\widehat{SE}=\sqrt{0.40(0.60)/100}\approx0.04899$. With 400 riders, $\sqrt{400}=2\sqrt{100}$, so the SE halves. Multiply each SE by 1.645 to obtain $E$ and use $\hat p\pm E$ for the endpoints.
	\end{minipage}

	\secInicio[prob:c7-1][sec:worked-solution-examples]

	\Needspace{5\baselineskip}
	\phantomsection
	\label{sol:c7-1-b}
	\subsubsection*{C7 -- Explain the gain in precision}

	The widths are approximately $2(8.06)=16.12$ and $2(4.03)=8.06$ percentage points. Quadrupling $n$ halves the width here because the observed share and the confidence level are the same. The larger sample locates the population share of mobile-ticket buyers more precisely, with both procedures targeting 90\% long-run coverage.

	\secInicio[prob:c7-1][sec:worked-solution-examples]

	\clearpage

	% ----------------------------------------------------------------
	\subsection{Mobile ticketing and the confidence target (C7)}
	\label{prob:c7-2}

	Use the rider surveys in Problem 1. The sample with 400 riders has $\hat p=0.40$ and $\widehat{SE}=\sqrt{0.40(0.60)/400}$. Its 90\% interval is about $[35.97\%,44.03\%]$. The 90\% interval for 100 riders is about $[31.94\%,48.06\%]$. The authority now considers 99\% confidence, with $z^*=2.576$.

	\QuestionsTasks[sec:worked-solution-examples]{
		{7}/{sol:c7-2-a}/
		{Construct the 99\% interval for 400 riders. Compare its margin, width and nominal coverage with the 90\% interval from those same 400 riders.},
		{7}/{sol:c7-2-b}/
		{Compare the 99\% interval from 400 riders with the 90\% interval from 100 riders. Explain why the result is compatible with the trade-off above. Which meets a requirement of at least 95\% confidence and a margin below 7 percentage points?}
	}

	\Needspace{5\baselineskip}
	\phantomsection
	\label{sol:c7-2-a}
	\subsubsection*{C7 -- Raise confidence while holding sample size fixed}

	The 99\% interval has the same centre of $0.40$ and the same SE of about 0.02449. Its larger critical value gives
	\[
	E=2.576\sqrt{\frac{0.40(0.60)}{400}}\approx0.06310,\qquad
	0.40\pm E\ \Longrightarrow\ [33.69\%,46.31\%]\text{ (rounded)}.
	\]
	The margin increases from about 4.03 to 6.31 percentage points, so the width rises from 8.06 to 12.62 percentage points. That loss of precision comes with higher nominal coverage, from 90\% to 99\%.

	\secInicio[prob:c7-2][sec:worked-solution-examples]

	\Needspace{5\baselineskip}
	\phantomsection
	\label{sol:c7-2-b}
	\subsubsection*{C7 -- Compare plans when both factors change}

	The 400-rider 99\% interval is narrower (about 12.62 versus 16.12 percentage points) and has higher nominal coverage. The larger sample halves the SE at this common observed share. That more than offsets the increase in the multiplier from 1.645 to 2.576. The 400-rider 99\% plan meets both requirements, with a 6.31 percentage point margin. The 100-rider 90\% plan misses both.

	\secInicio[prob:c7-2][sec:worked-solution-examples]


\end{document}
